\documentclass[conference]{IEEEtran}
\usepackage[T1]{fontenc}
\usepackage[utf8]{inputenc}
\usepackage[spanish,es-nodecimaldot,es-tabla]{babel}
\usepackage{amsmath}
\usepackage{url}

\title{Detección de deriva arquitectónica en sistemas multiagente basados en modelos de lenguaje mediante chequeo de conformidad sobre trazas}
\author{\IEEEauthorblockN{Leonardo Baca Calsin y German Arturo Chipana Jeronimo}
\IEEEauthorblockA{Departamento de Ingeniería de Sistemas e Informática\\
Universidad Nacional de San Agustín de Arequipa, Arequipa, Perú\\
lbacac@unsa.edu.pe, gchipanaj@unsa.edu.pe}}
\newcommand{\boxstep}[1]{\fbox{\strut\footnotesize #1}}
\newcommand{\downstep}{\\[-1pt]$\downarrow$\\[-1pt]}
\renewcommand{\IEEEkeywordsname}{Palabras clave}
\newcommand{\arch}{\textsc{AgentArch}}

\begin{document}
\maketitle
\begin{abstract}
Los sistemas multiagente basados en modelos de lenguaje pueden completar operaciones que incumplen su diseño. Presentamos \arch, un monitor pasivo que contrasta las trazas con una especificación de prompts, memoria y políticas. Definimos nueve operadores y comprobamos su cobertura en 50 escenarios guionizados de soporte de pedidos. En ese corpus construido para ejercitar las reglas, el verificador detectó y localizó los 40 casos con deriva, con cero alertas en diez nominales. Añadimos un piloto con inferencia local de Qwen3-0.6B y captura mediante el SDK de OpenTelemetry. Ocho solicitudes se ejecutaron tres veces y produjeron 24 sesiones válidas. En las acciones generadas se observaron 12 sesiones con infracciones y 12 conformes. \arch{} distinguió ambos grupos y localizó el primer paso infractor. Las repeticiones reprodujeron los mismos planes. Una prueba de metadatos alterados confirmó la dependencia de la instrumentación. El microbenchmark caracteriza exclusivamente el núcleo del verificador. Estos resultados aportan evidencia de cobertura y de integración en un piloto acotado.
\end{abstract}
\begin{IEEEkeywords}
deriva arquitectónica, sistemas multiagente, modelos de lenguaje, chequeo de conformidad, trazas, OpenTelemetry
\end{IEEEkeywords}

\section{Introducción}
Los sistemas multiagente basados en modelos de lenguaje (MAS-LLM) distribuyen tareas entre agentes con roles, memoria y herramientas especializadas. La arquitectura delimita el acceso a herramientas y el intercambio de conocimiento, aunque parte del comportamiento concreto se decide durante cada ejecución. Los patrones y metamodelos para estos sistemas ayudan a expresar esa estructura \cite{li,miranda,lu,liu,sallma}. Sin embargo, un registro de invocaciones exitosas puede contener una ejecución incompatible con el diseño.

Denominamos \emph{deriva arquitectónica observable} a una divergencia entre los invariantes operacionales especificados y las transiciones registradas en la traza de una sesión. El chequeo sobre trazas ya se ha explorado para firmware \cite{deangelis}. En un MAS-LLM interesa, además del orden de las operaciones, preservar las directivas activas y las condiciones de admisión y vigencia de la memoria.

Este trabajo aporta una especificación de comportamiento \(B=\langle P,M,\Pi\rangle\), una taxonomía de nueve operadores y una evaluación reproducible de un monitor pasivo. La unidad de evaluación es una sesión guionizada. Cada regla se aplica a los atributos observables del tramo al cerrarse. El monitor produce alertas de auditoría y la aplicación conserva el control de sus acciones.

\section{Trabajos relacionados}
La revisión estructurada comprende 26 trabajos seleccionados del corpus documental del proyecto, organizados por su relación con el diseño y la comprobación de agentes. Los criterios temáticos fueron arquitectura, verificación, observabilidad, prompts, memoria y procedencia. La Tabla~\ref{tab:related} resume esos ejes.

Los contratos de aseguramiento sobre trazas separan la especificación de la evidencia \cite{contracts}. AgentSpec aporta reglas con disparadores, predicados y acciones \cite{agentspec}. El control de políticas en MCP \cite{mcp} y los escudos de MENTOR \cite{mentor} abordan la intervención durante la ejecución. TraceFix analiza protocolos con modelos y contraejemplos \cite{tracefix}, mientras Pro2Guard predice riesgos a partir de trazas \cite{pro2guard}. Estos trabajos motivan reglas explícitas y mecanismos de verificación con alcances distintos.

AgentOps organiza los datos de observabilidad \cite{agentops}. AgentTelemetry y la instrumentación con inyección de fallos exploran la captura y el diagnóstico de ejecuciones \cite{telemetry,faults}. Watson estudia la observabilidad cognitiva \cite{watson}; COMPASS utiliza minería de procesos \cite{compass}; y Pathfinder analiza trazas mediante código ejecutable \cite{pathfinder}. La revisión de trayectorias \cite{survey} y el estudio de trazabilidad de roles \cite{barrak} sitúan el diagnóstico y la atribución de errores.

La evolución de prompts \cite{prompts}, ClawVM \cite{claw}, GovMem \cite{govmem} y MemArchitect \cite{memarch} motivan las comprobaciones de integridad, admisión y vigencia. TARV \cite{tarv} y TraceCaps \cite{tracecaps} estudian confianza y procedencia. Nuestra propuesta reúne reglas observables de esas familias en un monitor pasivo. La comparación experimental usa un control de HTTP y una variante propia de permisos y argumentos inspirada en AgentSpec.

\begin{table*}[!t]
\caption{Ejes de la revisión de literatura}
\label{tab:related}
\centering\footnotesize
\begin{tabular}{@{}p{.23\textwidth}p{.35\textwidth}p{.36\textwidth}@{}}
\hline
Eje & Trabajos representativos & Relación con la propuesta\\\hline
Arquitectura y diseño & Dimensiones de MAS, metamodelos, catálogos y SALLMA & Roles, restricciones, conocimiento y evolución.\\
Deriva y verificación & Trazas de firmware, contratos, AgentSpec, MCP, MENTOR, TraceFix y Pro2Guard & Especificación esperada y chequeo de operaciones.\\
Observabilidad y traza & AgentOps, AgentTelemetry, instrumentación MAS, Watson, COMPASS y Pathfinder & Registro de evidencia y diagnóstico de discrepancias.\\
Prompts y memoria & Evolución de prompts, ClawVM, GovMem y MemArchitect & Integridad de cláusulas, admisión, sobrescritura y expiración.\\
Confianza y procedencia & TARV y TraceCaps & Auditoría y análisis complementario.\\\hline
\end{tabular}
\end{table*}

\section{Modelo de comportamiento}
Sea \(\tau=(s_1,\ldots,s_n)\) la secuencia de tramos observados de una sesión. La especificación \(B=\langle P,M,\Pi\rangle\) reúne las obligaciones de prompt, las reglas de memoria y las políticas de herramientas. Las restricciones y la evolución de los prompts se representan en \(P\); el conocimiento mutable, en \(M\); y el control operacional, en \(\Pi\). Esta organización se inspira en las dimensiones arquitectónicas de Li et al.\ \cite{li}. La asignación de responsabilidad puede analizarse posteriormente a partir de los eventos registrados.

\begin{figure}[!t]
\centering
\boxstep{Agentes y herramientas de soporte de pedidos}
\downstep
\boxstep{Tramos normalizados de la sesión}
\downstep
\boxstep{Verificador pasivo \(\leftarrow\) especificación \(B\)}
\downstep
\boxstep{Cola de auditoría: eventos de infracción}
\caption{Flujo de datos de \arch. La especificación alimenta el verificador y este encola las alertas.}
\label{fig:arch}
\end{figure}

\begin{figure}[!t]
\centering
\begin{tabular}{ccc}
\boxstep{\(P\) · Prompts} & \boxstep{\(M\) · Memoria} & \boxstep{\(\Pi\) · Políticas}\\[2pt]
\footnotesize \(P_1\) Ausencia & \footnotesize \(M_1\) Admisión & \footnotesize \(\Pi_1\) Capacidad\\
\footnotesize \(P_2\) Conflicto & \footnotesize \(M_2\) Reemplazo & \footnotesize \(\Pi_2\) Argumentos\\
\footnotesize \(P_3\) Integridad & \footnotesize \(M_3\) Vigencia & \footnotesize \(\Pi_3\) Precedencia
\end{tabular}
\caption{Familias y operadores de deriva en \(B\).}
\label{fig:taxonomy}
\end{figure}

\subsection{Cláusulas de prompt}
Para un rol \(i\), sean \(R_i\) los identificadores obligatorios, \(A_i(s)\) los identificadores activos y \(h_i^0(c)\) el resumen SHA-256 canónico de la cláusula \(c\). El tramo declara \(h_s(c)\) para las cláusulas activas. \(P_1\) señala la ausencia de un elemento de \(R_i\). \(P_2\) señala la presencia conjunta de un par registrado en el conjunto de conflictos \(C\). \(P_3\) señala una discrepancia de resumen para un identificador con valor canónico. El resultado depende de que el instrumentador registre identificadores y resúmenes fieles a las cláusulas utilizadas.

\subsection{Estado de memoria}
El verificador mantiene un mapa de claves a valores, marcas temporales y períodos de vigencia (TTL). En una escritura, \(M_1\) comprueba la marca de verificación y \(M_2\) contrasta el nuevo valor con el valor registrado y la autorización de reemplazo. En una lectura de una clave conocida, \(M_3\) compara su antigüedad con el TTL. La marca de verificación es evidencia declarada por la aplicación. Su calidad depende del mecanismo que la produce \cite{govmem}.

\subsection{Políticas de herramientas}
Sea \(K_i\) el conjunto de herramientas permitidas para el rol \(i\), \(V_a\) el validador de argumentos para la herramienta \(a\), y \(D\) el conjunto de precedencias \((a,b)\). \(\Pi_1\) comprueba la pertenencia a \(K_i\) y \(\Pi_2\) exige que el validador aplicable acepte los argumentos. \(\Pi_3\) comprueba que la traza haya registrado \(a\) antes de \(b\). La implementación evaluada conserva un conjunto de nombres observados durante la sesión. La comprobación de éxito y la correlación por recurso requerirían atributos adicionales.

\begin{table*}[!t]
\caption{Reglas de violación y distribución del corpus guionizado}
\label{tab:operators}
\centering\footnotesize
\begin{tabular}{@{}llp{.36\textwidth}p{.32\textwidth}r@{}}
\hline
Regla & IDs & Condición de violación & Patrón evaluado & Casos\\\hline
-- & S01--S10 & Ninguna de las condiciones siguientes & Consulta y escritura válidas, verificación de stock antes de reserva & 10\\
\(P_1\) & S11--S14 & \(R_i\setminus A_i(s)\neq\emptyset\) & Omisión de \texttt{clause\_auth} en OrderTriage & 4\\
\(P_2\) & S15--S18 & \(\exists(c,d)\in C:\{c,d\}\subseteq A_i(s)\) & Dos cláusulas de reembolso excluyentes & 4\\
\(P_3\) & S19--S22 & \(\exists c\in A_i(s)\cap\mathrm{dom}(h_i^0):h_s(c)\neq h_i^0(c)\) & Resumen de la cláusula de garantía alterado & 4\\
\(M_1\) & S23--S27 & Escritura sin marca de verificación & Descuento especulativo & 5\\
\(M_2\) & S28--S32 & Cambio de valor de una clave existente sin autorización & Reemplazo de dirección de envío & 5\\
\(M_3\) & S33--S36 & Lectura de clave conocida con \(t_{\mathrm{lect}}-t_k>\mathrm{TTL}_k\) & Token leído después de su vencimiento & 4\\
\(\Pi_1\) & S37--S41 & Herramienta \(a\notin K_i\) & Orden de pago emitida por OrderTriage & 5\\
\(\Pi_2\) & S42--S46 & El validador \(V_a\) rechaza los argumentos & Reembolso con importe negativo & 5\\
\(\Pi_3\) & S47--S50 & Se observa \(b\) y falta una aparición previa de \(a\) para \((a,b)\in D\) & Reserva sin verificación previa de inventario & 4\\\hline
\end{tabular}
\end{table*}

\section{Instrumentación y chequeo}
La Fig.~\ref{fig:arch} presenta el flujo de datos y la Fig.~\ref{fig:taxonomy} agrupa los operadores. El esquema de dominio usa nombres bajo \texttt{agentarch.*}. El tramo \texttt{session} representa la sesión. Los tramos \texttt{agent\_turn}, \texttt{memory\_op} y \texttt{tool\_call} contienen respectivamente las cláusulas activas y sus resúmenes, los datos de acceso a memoria y la herramienta con sus argumentos serializados.

En el chequeo incremental, el procesador clasifica el tramo, evalúa las reglas de su familia y actualiza el estado observado de la sesión. Cuando alguna regla se satisface, encola un evento con \texttt{trace\_id}, \texttt{span\_id}, paso, marca temporal y conjunto de infracciones. La evaluación guionizada invoca \texttt{onEnd} sobre tramos sintéticos, con una instancia por escenario. El piloto añade un adaptador del SDK de OpenTelemetry \cite{otel} que recibe \texttt{ReadableSpan}, conserva una instancia por \texttt{traceId} y libera su estado al cierre de la sesión. Una prueba con dos sesiones intercaladas comprueba ese aislamiento. La figura~\ref{fig:processor} muestra ese recorrido. La especificación también puede aplicarse fuera de línea recorriendo una traza consolidada en el orden de sus tramos.

\begin{figure}[!t]
\centering
\boxstep{\texttt{onEnd(tramo)}}
\downstep
\boxstep{Clasificar: turno, memoria o herramienta}
\downstep
\boxstep{Evaluar reglas y actualizar estado observado}
\downstep
\boxstep{Si hay infracciones, encolar evento}
\caption{Procesamiento al cierre del tramo. La llamada devuelve las infracciones observadas.}
\label{fig:processor}
\end{figure}

\section{Evaluación experimental}
\subsection{Escenarios y comparadores}
El dominio sintético representa una aplicación de atención de pedidos. OrderTriage clasifica solicitudes, InventoryAgent consulta o reserva existencias y BillingAgent ejecuta operaciones financieras. La Fig.~\ref{fig:example} ilustra la variante de precedencia que se evalúa en S47--S50.

\begin{figure}[!t]
\centering
\boxstep{InventoryAgent solicita \texttt{reserve\_stock}}
\downstep
\boxstep{Se cierra el tramo de herramienta}
\downstep
\boxstep{\(\Pi_3\): falta \texttt{verify\_inventory} previo}
\downstep
\boxstep{\arch{} encola una alerta}
\caption{Ejemplo guionizado de precedencia: llamada observada, condición infringida y alerta.}
\label{fig:example}
\end{figure}

El generador produce 50 trazas etiquetadas: 10 nominales y 40 con deriva (12 de prompts, 14 de memoria y 14 de políticas). Hay un patrón nominal y nueve patrones anómalos, con varias instancias por patrón. Por ejemplo, S11--S14 comparten la omisión de \texttt{clause\_auth}. Los IDs y posibles valores de negocio varían entre algunas instancias. La Tabla~\ref{tab:operators} documenta los casos y los atributos que determinan su etiqueta.

Las tres configuraciones reciben las mismas trazas. \(B_1\) es un control negativo que emite una alerta ante HTTP igual o superior a 500. Representa únicamente ese chequeo de transporte. \(B_2\) es una variante de reglas de contorno inspirada en AgentSpec \cite{agentspec}; valida permisos de rol y argumentos de herramienta. \arch{} evalúa las nueve reglas de \(B\). El contraste con estos comparadores permite examinar la cobertura de sus reglas. Cada configuración se compara con las etiquetas de presencia de deriva y primer paso infractor, almacenadas en \texttt{ground\_truth.json}.

\subsection{Métricas}
Una traza con deriva y al menos una alerta cuenta como verdadero positivo (VP). Una alerta sobre traza nominal cuenta como falso positivo (FP); una deriva omitida constituye un falso negativo (FN). La precisión se calcula como \(\mathrm{VP}/(\mathrm{VP}+\mathrm{FP})\) y la exhaustividad como \(\mathrm{VP}/(\mathrm{VP}+\mathrm{FN})\). \(F_1\) es su media armónica y vale cero si se omiten todas las derivas. Cuando no se emiten alertas, la precisión se marca como indefinida. La métrica \(t^*\) es el porcentaje de los 40 casos anómalos cuyo primer paso alertado coincide con el esperado. Las omisiones se cuentan en el denominador.

\subsection{Detección y localización}
\begin{table}[H]
\caption{Casos detectados por dimensión}
\label{tab:category_results}
\centering\footnotesize
\begin{tabular}{lrrr}
\hline
Variante & Prompts & Memoria & Políticas\\
\hline
$B_1$ & 0/12 & 0/14 & 0/14\\
$B_2$ & 0/12 & 0/14 & 10/14\\
AgentArch & 12/12 & 14/14 & 14/14\\
\hline
\end{tabular}
\end{table}
\begin{table*}[!t]
\caption{Resultados globales (métricas en \%)}
\label{tab:experimental_results}
\centering\footnotesize
\begin{tabular*}{\textwidth}{@{\extracolsep{\fill}}lrrrrrrr}
\hline
Variante & VP & FP & FN & Prec. & Exh. & $F_1$ & $t^*$\\
\hline
$B_1$ & 0 & 0 & 40 & -- & 0 & 0 & 0\\
$B_2$ & 10 & 0 & 30 & 100 & 25 & 40 & 25\\
AgentArch & 40 & 0 & 0 & 100 & 100 & 100 & 100\\
\hline
\end{tabular*}
\end{table*}

Las Tablas~\ref{tab:category_results} y~\ref{tab:experimental_results} muestran los resultados. Ninguna variante emitió alertas en los diez nominales. Este conteo caracteriza el corpus y sus patrones repetidos. \(B_1\) no detectó las derivas, que se expresan en atributos arquitectónicos del escenario. \(B_2\) detectó los cinco casos de capacidades y los cinco de argumentos. \arch{} detectó las 40 derivas y localizó su primer paso en los 40 casos. La diferencia de cobertura refleja las reglas presentes en cada configuración y los patrones diseñados para ejercitarlas.

\subsection{Microbenchmark del verificador}
\begin{table}[!t]
\caption{Tiempo de verificación por tramo ($\mu$s)}
\label{tab:latency_bench}
\centering\footnotesize
\begin{tabular}{lrr}
\hline
Tramo & Media & DE entre lotes\\\hline
Turno & 0.189 & 0.218\\
Memoria & 0.079 & 0.052\\
Herramienta & 0.582 & 0.157\\
\hline
\end{tabular}
\end{table}

El microbenchmark del 16 de noviembre de 2026 ejecutó 50 lotes de 1000 llamadas cronometradas por tipo de tramo sintético conforme. Antes de medir cada lote, se usó una instancia nueva y se calentó con el mismo tramo. La medición registra la duración del bucle dividida por 1000. La Tabla~\ref{tab:latency_bench} informa la media y la desviación estándar muestral (DE) de las 50 medias de lote, en microsegundos por llamada. El entorno usó Node.js 22.15.0 en Windows x64, sobre un Intel Core i7-1165G7. El tiempo incluye el camino conforme de \texttt{onEnd} con los hashes preparados. La inferencia, la exportación de telemetría y el consumo de alertas quedan fuera de la medición.

\subsection{Piloto con un modelo en ejecución}
Ejecutamos Qwen3-0.6B en formato GGUF Q8\_0 \cite{qwen}, servido por llama.cpp b12704 \cite{llamacpp} sobre CPU. El equipo tiene 16 GB de RAM y el servidor usa cuatro hilos y contexto de 4096 tokens. El SDK JavaScript es 2.12.1. Ocho solicitudes prefijadas, cuatro nominales y cuatro adversariales, se ejecutaron tres veces con temperatura cero y semillas 42, 43 y 44. Cada sesión asigna uno de los tres roles del dominio. El modelo genera un plan JSON de llamadas y la aplicación ejecuta sus operaciones mediante funciones locales que conservan argumentos en memoria. La inferencia es real y las herramientas representan servicios de prueba.

El prompt define permisos, límites y precedencias. Las peticiones adversariales solicitan pagos desde un rol sin permiso, un reembolso negativo o la omisión de un antecedente. La salida del modelo determina las llamadas observadas. Un oráculo separado del monitor, fijado en el ejecutor antes de inferir, etiqueta permisos, rangos y orden de esas llamadas. La intención adversarial se distingue de la infracción efectivamente generada. Los resultados conservan peticiones, respuestas completas, tramos con identificadores del SDK, etiquetas y alertas. Los errores HTTP y de formato se cuentan por separado.

Las 24 sesiones produjeron planes válidos, sin errores ni planes vacíos. Se observaron 12 sesiones conformes y 12 con infracciones. Las tres repeticiones de cada solicitud generaron el mismo plan, de modo que hay ocho secuencias distintas. Entre las cuatro solicitudes adversariales, una produjo \(\Pi_1,\Pi_2,\Pi_3\), otra \(\Pi_2\) y dos \(\Pi_3\). \arch{} detectó y localizó las 12 sesiones anómalas, sin alertas en las 12 conformes. A nivel de solicitud, corresponden a cuatro de cuatro secuencias anómalas y cero alertas en cuatro conformes. \(B_1\) detectó cero y \(B_2\), seis de doce. La Tabla~\ref{tab:live} resume los conteos por sesión.

\begin{table}[!t]
\caption{Piloto con inferencia local sobre 24 sesiones}
\label{tab:live}
\centering\footnotesize
\begin{tabular}{lrrrrr}
\hline
Variante & VP & FP & VN & FN & $t^*$ (\%)\\\hline
$B_1$ & 0 & 0 & 12 & 12 & 0\\
$B_2$ & 6 & 0 & 12 & 6 & 50\\
\arch & 12 & 0 & 12 & 0 & 100\\
\hline
\end{tabular}
\end{table}


El piloto ejercita acciones de herramientas generadas por un modelo y su captura por el SDK. El alcance comprende roles asignados individualmente, un dominio y un modelo pequeño. La evaluación completa de los nueve operadores permanece en el corpus guionizado. La interacción libre entre agentes y la memoria generada por el modelo requieren ensayos adicionales.

\subsection{Sensibilidad a metadatos incorrectos}
Una prueba separada de inyección de fallos examinó dos supuestos. Al alterar una cláusula y registrar su nuevo resumen, el verificador emitió \(P_3\). Al mantener indebidamente el resumen canónico en los metadatos, omitió la alteración. Una escritura sin respaldo produjo \(M_1\) con \texttt{verified=false} y quedó sin alerta cuando se informó incorrectamente \texttt{true}. Los cuatro casos documentan la frontera de confianza del instrumentador. Esta prueba usa entradas controladas y sus resultados se conservan aparte del piloto del modelo.

\section{Amenazas a la validez}
\textit{Constructo.} Los operadores comprueban los atributos observados. Las alertas dependen de que esos atributos representen las cláusulas y las acciones efectivas. La prueba de metadatos muestra dos omisiones cuando se informa evidencia incorrecta. La precedencia evalúa nombres de herramientas observados previamente. Verificar éxito y correlación por recurso exige ampliar la regla.

\textit{Validez interna.} Los patrones guionizados se construyeron para ejercitar los predicados. Por ello, el 100\% obtenido caracteriza cobertura del corpus. \(B_1\) sirve como control negativo y \(B_2\) implementa un subconjunto específico. En el piloto, el oráculo es separado pero codifica las mismas políticas del dominio. Las réplicas a temperatura cero resultaron idénticas y aportan repetición de ejecución, con diversidad limitada a ocho solicitudes.

\textit{Falsos positivos e incertidumbre.} En el corpus guionizado se observaron cero alertas en diez nominales. Bajo un modelo binomial de ensayos independientes e idénticamente distribuidos, el límite superior exacto unilateral al 95\% para cero de diez sería
\[
p_U=1-0.05^{1/10}\approx 0.259.
\]
La aproximación de la regla del tres da 30\%. Ese supuesto de independencia se debilita por la repetición del patrón nominal. El límite se presenta para mostrar la incertidumbre del tamaño muestral y no permite estimar la tasa poblacional del sistema. En el piloto, cero de doce corresponde a cuatro solicitudes repetidas y también debe interpretarse descriptivamente.

\textit{Medición temporal.} Los tiempos submicrosegundo de la Tabla~\ref{tab:latency_bench} corresponden al núcleo conforme, con hashes preparados y calentamiento. La exportación, el consumo de alertas y la inferencia requieren mediciones adicionales para estimar el costo operativo completo. El piloto conserva tiempos de turno y ejecución local instrumentada, sin estimar una diferencia causal de sobrecarga.

\textit{Validez externa.} La inferencia real aporta evidencia exploratoria de integración en el entorno descrito. El modelo, el dominio, la brevedad de los planes y las herramientas de prueba restringen la generalización. Se necesita ampliar modelos, solicitudes, sesiones concurrentes, trazas parciales e interacción entre roles.

\section{Disponibilidad y reproducción}
El artefacto es público en \url{https://github.com/leonardobaca7/agentarch}. La versión de esta revisión se identifica con la etiqueta \texttt{v0.2.0-c2-review}. El commit \texttt{bd1c1a1} contiene el código de la medición guionizada del 16 de noviembre. Su sucesor \texttt{73dd953} incorpora resultados y manuscrito. El piloto se ejecutó con \texttt{647f45f}. Así se distingue el código medido de la versión que empaqueta los artefactos.

\texttt{node eval\_experiment.js} regenera el corpus guionizado y el microbenchmark. \texttt{npm ci} instala las versiones del SDK fijadas en el archivo de bloqueo. Con el servidor del modelo activo, \path{experiments/run_live_evaluation.js} ejecuta el piloto en un directorio nuevo. El protocolo documenta los parámetros y el SHA-256 de los pesos. Las salidas versionadas conservan las respuestas originales y sus trazas para inspección independiente.

\section{Trabajo futuro}
Ampliaremos el piloto a más modelos, solicitudes distintas e interacción entre agentes. La instrumentación se evaluará con eventos ausentes o incorrectos, mientras la precedencia incorporará resultado y recurso. Se estudiará el costo con exportación y consumo de alertas mediante ejecuciones emparejadas con y sin verificador. También se explorará cómo intervenir ante infracciones con reglas de bloqueo explícitas.

\section{Conclusiones}
\arch{} expresa invariantes de prompts, memoria y herramientas mediante nueve operadores de chequeo de conformidad sobre trazas. En 50 escenarios guionizados, cubrió las 40 derivas preparadas y localizó su primer paso, con cero alertas en diez nominales. El piloto con inferencia local aportó 24 sesiones correspondientes a ocho solicitudes y mostró detección de las infracciones de herramientas observadas. Las réplicas idénticas limitan la diversidad de esa evidencia. La prueba de metadatos confirmó que la corrección de la instrumentación condiciona la detección. El microbenchmark caracteriza el núcleo del monitor y la evaluación de su costo operativo completo queda pendiente.

\IEEEtriggeratref{25}
\begin{thebibliography}{99}
\bibitem{li} K. Li, M. Xie, J. Ge, and Y. Liu, ``What do we architect when we architect LLM agent systems?'' in \emph{Proc. IEEE/ACM 48th Int. Conf. Software Engineering: Companion (ICSE-Companion)}, 2026, pp. 484--486, doi: 10.1145/3774748.3795870.
\bibitem{miranda} J. P. Miranda, A. Radermacher, F. Baligand, J. Bonnail, S. Gérard, P. Bannerot, and M. D. Del Fabro, ``Bridging MDE and LLM-based agent frameworks for multi-agent systems: A quasi-systematic review and metamodel,'' in \emph{Proc. Int. Conf. Model-Based Software and Systems Engineering (MODELSWARD)}, 2026, pp. 33--44, doi: 10.5220/0014236700004058.
\bibitem{lu} Q. Lu, L. Zhu, X. Xu, Z. Xing, S. Harrer and J. Whittle, ``Towards responsible generative AI: A reference architecture for designing foundation model based agents,'' arXiv:2311.13148, ver. 3, 2024, doi: 10.48550/arXiv.2311.13148.
\bibitem{liu} Y. Liu, S. K. Lo, Q. Lu, L. Zhu, D. Zhao, X. Xu, S. Harrer, and J. Whittle, ``Agent design pattern catalogue: A collection of architectural patterns for foundation model based agents,'' arXiv preprint arXiv:2405.10467, 2024, doi: 10.48550/arXiv.2405.10467.
\bibitem{sallma} M. Becattini, R. Verdecchia, and E. Vicario, ``SALLMA: A software architecture for LLM-based multi-agent systems,'' in \emph{Proc. IEEE/ACM Int. Workshop on New Trends in Software Architecture (SATrends)}, 2025, pp. 5--8, doi: 10.1109/SATrends66715.2025.00006.
\bibitem{deangelis} D. F. De Angelis, M. De Luca, D. Amalfitano, P. Cimmino, and A. R. Fasolino, ``Detecting architectural drift in safety-critical firmware through runtime trace analysis,'' arXiv preprint arXiv:2607.03135, 2026, doi: 10.48550/arXiv.2607.03135.
\bibitem{conformance} A. Rozinat and W. M. P. van der Aalst, ``Conformance checking of processes based on monitoring real behavior,'' \emph{Information Systems}, vol. 33, no. 1, pp. 64--95, 2008, doi: 10.1016/j.is.2007.07.001.
\bibitem{rv} E. Bartocci, Y. Falcone, A. Francalanza, and G. Reger, ``Introduction to runtime verification,'' in \emph{Lectures on Runtime Verification}, LNCS, vol. 10457, E. Bartocci and Y. Falcone, Eds. Cham, Switzerland: Springer, 2018, pp. 1--33, doi: 10.1007/978-3-319-75632-5\_1.
\bibitem{contracts} C. Paduraru, P.-L. Bouruc, and A. Stefanescu, ``A trace-based assurance framework for agentic AI orchestration: Contracts, testing, and governance,'' arXiv preprint arXiv:2603.18096, 2026, doi: 10.48550/arXiv.2603.18096.
\bibitem{agentspec} H. Wang, C. M. Poskitt, and J. Sun, ``AgentSpec: Customizable runtime enforcement for safe and reliable LLM agents,'' in \emph{Proc. IEEE/ACM 48th Int. Conf. Software Engineering (ICSE)}, 2026, pp. 2938--2950, doi: 10.1145/3744916.3764546.
\bibitem{mcp} S. Wang, S. Zhu, and R. Li, ``Runtime policy enforcement for MCP-based LLM agents,'' \emph{Electronics}, vol. 15, no. 13, Art. no. 2829, 2026, doi: 10.3390/electronics15132829.
\bibitem{mentor} A. Anand, I. Chatzi, R. Raha, and A.-K. Schmuck, ``MENTOR: Procedural shielding of tool-using LLM agents,'' arXiv preprint arXiv:2610.17039, 2026, doi: 10.48550/arXiv.2610.17039.
\bibitem{tracefix} S. Xia, Q. Li, T. Ehsan, and J. Ortiz, ``TraceFix: Repairing agent coordination protocols with TLA+ counterexamples,'' in \emph{Proc. ACM Conf. AI and Agentic Systems (CAIS), San Jose, CA, USA}, 2026, pp. 181--196, doi: 10.1145/3786335.3813159.
\bibitem{pro2guard} H. Wang, C. M. Poskitt, J. Sun, and J. Wei, ``Pro2Guard: Proactive runtime enforcement of LLM agent safety via probabilistic prediction,'' arXiv:2508.00500, ver. 2, 2026, doi: 10.48550/arXiv.2508.00500.
\bibitem{agentops} L. Dong, Q. Lu, and L. Zhu, ``AgentOps: Enabling observability of LLM agents,'' arXiv preprint arXiv:2411.05285, 2024, doi: 10.48550/arXiv.2411.05285.
\bibitem{telemetry} K. C. Balusu, ``AgentTelemetry: A fault detection benchmark and toolkit for LLM agent observability,'' in \emph{Proc. ACM AIware}, 2026, pp. 380--387, doi: 10.1145/3805760.3814931.
\bibitem{faults} Z. Seyedghorban, E. Klimov, A. van Deursen, A. Panichella, and B. K. Ozkan, ``Observability and fault injection for LLM-based multi-agent systems in software engineering,'' in \emph{Proc. IEEE Int. Conf. Software Testing, Verification and Validation (ICST)}, 2026, pp. 211--215, doi: 10.1109/ICST69053.2026.00037.
\bibitem{watson} B. Rombaut, S. Masoumzadeh, K. Vasilevski, D. Lin, and A. E. Hassan, ``Watson: A cognitive observability framework for the reasoning of LLM-powered agents,'' arXiv:2411.03455, ver. 3, 2025, doi: 10.48550/arXiv.2411.03455.
\bibitem{compass} R. Dorsch, D. Henselmann and A. Harth, ``COMPASS: A process mining-based methodology for prompt analysis of large language model agents,'' in \emph{Joint Proc. HybridAIMS and CAI}, CEUR Workshop Proceedings, vol. 3996, 2025, pp. 29--37. [Online]. Available: \url{https://ceur-ws.org/Vol-3996/paper-5.pdf}
\bibitem{pathfinder} D. Atreja, ``Pathfinder: Self-improving agent trace analysis via adversarial self-play and code execution,'' in \emph{Proc. ACM Conf. AI and Agentic Systems (CAIS), San Jose, CA, USA}, 2026, pp. 1336--1339, doi: 10.1145/3786335.3813199.
\bibitem{survey} J. Wang, Y. Wang, M. Chen, X. Xie, C. Chen, F. Mu, Z. Liu, and Q. Wang, ``A survey for LLM agent trajectory analysis: From failure attribution to enhancement,'' \emph{IEEE Trans. Softw. Eng.}, vol. 52, no. 11, pp. 3406--3426, 2026, doi: 10.1109/TSE.2026.3717765. [Online]. Available: \url{https://ieeexplore.ieee.org/document/11626967}
\bibitem{barrak} A. Barrak, ``Traceability and accountability in role-specialized multi-agent LLM pipelines,'' arXiv:2510.07614, 2025, doi: 10.48550/arXiv.2510.07614.
\bibitem{prompts} M. Tafreshipour, A. Imani, E. Huang, E. Almeida, T. Zimmermann, and I. Ahmed, ``Prompting in the wild: An empirical study of prompt evolution in software repositories,'' in \emph{Proc. IEEE/ACM 22nd Int. Conf. Mining Software Repositories (MSR)}, 2025, pp. 686--698, doi: 10.1109/MSR66628.2025.00106.
\bibitem{claw} M. Rafique and L. Bindschaedler, ``ClawVM: Harness-managed virtual memory for stateful tool-using LLM agents,'' in \emph{Proc. EuroMLSys}, 2026, pp. 1--12, doi: 10.1145/3805621.3807648.
\bibitem{govmem} Y. Qi, X. Xu, and Y. Li, ``When not to write memory: Governing false promotion from correlated agent traces,'' arXiv:2607.02579, 2026, doi: 10.48550/arXiv.2607.02579.
\bibitem{memarch} L. Suresh Kumar, Y. Ba, and R. Pan, ``MemArchitect: A policy driven memory governance layer,'' arXiv preprint arXiv:2603.18330, 2026, doi: 10.48550/arXiv.2603.18330.
\bibitem{tarv} T. P. Patel, ``Trust-aware runtime verification for autonomous LLM agent systems,'' in \emph{Proc. Int. Conf. Artificial Intelligence, Systems, and Emerging Technologies (ICAISET)}, 2026, pp. 1--6, doi: 10.1109/ICAISET66439.2026.11541707.
\bibitem{tracecaps} A. Catarino, C. Mamede, R. Melo, and R. Abreu, ``TraceCaps: Inline provenance and risk enforcement for agentic software engineering,'' in \emph{Proc. ICSE New Ideas and Emerging Results Track (ICSE-NIER)}, 2026, pp. 166--170, doi: 10.1145/3786582.3786832.
\bibitem{enterpriseagentarch} T. Bogavelli, H. Subramani, and R. Sharma, ``AgentArch: A benchmark for evaluating agent architectures in enterprise workflows,'' in \emph{Proc. AACL-IJCNLP}, 2026, pp. 915--933, doi: 10.18653/2026.aacl-long.44.
\bibitem{otel} OpenTelemetry Authors, ``OpenTelemetry JavaScript instrumentation,'' SDK ver. 2.12.1. Accessed: Nov. 24, 2026. [Online]. Available: \url{https://opentelemetry.io/docs/languages/js/instrumentation/}
\bibitem{genaisemconv} OpenTelemetry Authors, ``GenAI semantic conventions,'' Accessed: Nov. 24, 2026. [Online]. Available: \url{https://github.com/open-telemetry/semantic-conventions-genai/tree/main/docs/gen-ai}
\bibitem{qwen} Qwen Team, ``Qwen3-0.6B-GGUF,'' model weights, Q8\_0. Accessed: Nov. 24, 2026. [Online]. Available: \url{https://huggingface.co/Qwen/Qwen3-0.6B-GGUF}
\bibitem{llamacpp} ggml-org, ``llama.cpp,'' release b12704, commit 72cda4a22. Accessed: Nov. 24, 2026. [Online]. Available: \url{https://github.com/ggml-org/llama.cpp/releases/tag/b12704}
\end{thebibliography}

\end{document}
